\scanpage{37}
% archive sheet number: 36
${\unsure{\mathscr{S}|S}}_\infty$ schéma d'anneaux \inserted{{\unsure{affine}}} sur $\mathbb{Z}$, si $A$ est un anneau, on a
\[
\mathscr{S}_\infty({\unsure{\mathbb{Z}|A}})=A[[t]]^{**}
\]
(séries formelles d'augmentation 1) où l'addition est la multiplication des séries formelles, et le produit est la convolution $*$. Il y a aussi des opérations $\lambda^i$. {\unsure{Rappelons}},
\[
(1+at)*f(t)=f(at)\qquad\text{i.e. }\{1+at\}\{f(t)\}=\{f(at)\}
\]
% margin note written sideways
\margin{{\unsure{Relations}} {\unsure{séries|liens|bien}} {\unsure{formelles}} et {\unsure{vecteurs de Witt}}}
Si $A$ est de car $p>0$, on aura un \uline{endomorphisme} $F$ de $\mathscr{S}_\infty(A)$ [provenant de l'endomorphisme de Frobenius de $\mathscr{S}_{\infty\,\mathbb{Z}/p\mathbb{Z}}$]
\[
F\left(\sum a_it^i\right)=\sum a_i^pt^i
\]
\figure{\begin{tikzpicture}\fill(0,0)circle(.015);\fill(.15,-.02)circle(.015);\draw(.25,.03)--(.4,.07);\draw(.48,-.12)--(.59,-.08);\fill(-.05,-.2)circle(.01);\end{tikzpicture}}
et on aura
\[
\deleted{{\unsure{F(f(at))}}}\quad f(t)^p=F(f)(t^p)
\]
Cela conduit alors à poser
\[
V(f)(t)=f(t^p)
\]
[$V$ est un {\unsure{endomorphisme|monomorphisme}} additif] et on aura
\[
FV(\{f\})=VF(\{f\})=p\{f\}.
\]
N.B. $F$ est aussi un endomorphisme pour la structure d'anneaux

En \deleted{{\unsure{a donc}}} par exemple
\[
(1+xt^{p^i})*f(t)=F^i(f)(xt^{p^i})
\]
car \deleted{{\unsure{si $x=a^{p^i}$, le premier membre est}} {\unsure{donc le membre}} ${\unsure{1+ax}}$}
\[
\begin{aligned}
(1+at)^{p^i}*f&=f(at)^{p^i}\\
&=F^i(f)(a^{p^i}t^{p^i})=F^i(f)(xt^{p^i}).
\end{aligned}
\]
Sans restriction \deleted{\illegible{1}} de caractéristique, on a un morphisme d'anneaux
\[
\psi:\mathscr{S}_\infty\longrightarrow W_\infty^{{\unsure{p}}}\qquad[\text{anneau de Witt pour le nb premier }p]
\]
\margin{?}
qui en tant que morphisme additif est déterminé par
\[
\psi(\{1+at\})=\epsilon(a)=(a,0,0,\ldots)
\]
